% Môn: Vật lý
% Lớp: 12
% Chương: Chương I. Vật lí nhiệt
% Bài: L12C1 Bài 6. Nhiệt hóa hơi riêng · Nhận biết khái niệm, đơn vị và ý nghĩa vật lí của nhiệt dung riêng
% ===== Câu 50 =====
% ID: L12C1B6-155-TL
% Mức: NB
\dangbt{Nhận biết khái niệm, đơn vị và ý nghĩa vật lí của nhiệt dung riêng}
\begin{bt}
% ID: L12C1B6-155-TL
% Mức: NB
Trình bày cách nhận biết và phương pháp giải dạng “Nhận biết khái niệm, đơn vị và ý nghĩa vật lí của nhiệt dung riêng”. Phân tích nhận định sau và sửa lại nếu sai: “Đơn vị SI của nhiệt dung riêng là J/kg.”
\loigiai{
Trước hết xác định hiện tượng, trạng thái và các đại lượng đã cho. Nêu định nghĩa hoặc hệ thức phù hợp, đổi về đơn vị SI, áp dụng đúng điều kiện và kiểm tra ý nghĩa kết quả. Nhận định cần đối chiếu: Nhiệt dung riêng cho biết nhiệt lượng cần để làm $1\,\text{kg}$ chất tăng 1 K. Vì vậy nhận định đã cho sai; phát biểu đúng là: Nhiệt dung riêng cho biết nhiệt lượng cần để làm $1\,\text{kg}$ chất tăng 1 K.
}
\end{bt}

% ===== Câu 54 =====
% ID: L12C1B6-156-TL
% Mức: NB
\begin{bt}
% ID: L12C1B6-156-TL
% Mức: NB
Trình bày cách nhận biết và phương pháp giải dạng “Nhận biết khái niệm, đơn vị và ý nghĩa vật lí của nhiệt dung riêng”. Phân tích nhận định sau và sửa lại nếu sai: “Đơn vị SI của nhiệt dung riêng là J/kg.”
\loigiai{
Trước hết xác định hiện tượng, trạng thái và các đại lượng đã cho. Nêu định nghĩa hoặc hệ thức phù hợp, đổi về đơn vị SI, áp dụng đúng điều kiện và kiểm tra ý nghĩa kết quả. Nhận định cần đối chiếu: Nhiệt dung riêng cho biết nhiệt lượng cần để làm $1\,\text{kg}$ chất tăng 1 K. Vì vậy nhận định đã cho sai; phát biểu đúng là: Nhiệt dung riêng cho biết nhiệt lượng cần để làm $1\,\text{kg}$ chất tăng 1 K.
}
\end{bt}

% ===== Câu 58 =====
% ID: L12C1B6-157-TL
% Mức: TH
\begin{bt}
% ID: L12C1B6-157-TL
% Mức: TH
Trình bày cách nhận biết và phương pháp giải dạng “Nhận biết khái niệm, đơn vị và ý nghĩa vật lí của nhiệt dung riêng”. Phân tích nhận định sau và sửa lại nếu sai: “Nhiệt dung riêng cho biết nhiệt lượng cần để làm $1\,\text{kg}$ chất tăng 1 K.”
\loigiai{
Trước hết xác định hiện tượng, trạng thái và các đại lượng đã cho. Nêu định nghĩa hoặc hệ thức phù hợp, đổi về đơn vị SI, áp dụng đúng điều kiện và kiểm tra ý nghĩa kết quả. Nhận định cần đối chiếu: Nhiệt dung riêng cho biết nhiệt lượng cần để làm $1\,\text{kg}$ chất tăng 1 K. Vì vậy nhận định đã cho đúng.
}
\end{bt}

% ===== Câu 61 =====
% ID: L12C1B6-158-DS
% Mức: NB
\begin{ex}
% ID: L12C1B6-158-DS
% Mức: NB
Xét các phát biểu sau về dạng “Nhận biết khái niệm, đơn vị và ý nghĩa vật lí của nhiệt dung riêng” (bộ số 4).
\choiceTF
{\True Nhiệt dung riêng cho biết nhiệt lượng cần để làm $1\,\text{kg}$ chất tăng 1 K.}
{Đơn vị SI của nhiệt dung riêng là J/kg.}
{\True Khi giải dạng “Nhận biết khái niệm, đơn vị và ý nghĩa vật lí của nhiệt dung riêng”, cần xác định đúng trạng thái đầu, trạng thái cuối và quy ước dấu hoặc điều kiện áp dụng.}
{Có thể bỏ qua mọi dữ kiện về khối lượng, nhiệt độ và hiệu suất trong mọi bài toán thuộc dạng này.}
\loigiai{
\begin{itemchoice}
\item (Đúng) Nhiệt dung riêng cho biết nhiệt lượng cần để làm $1\,\text{kg}$ chất tăng 1 K. (Vì): Nhiệt dung riêng cho biết nhiệt lượng cần để làm $1\,\text{kg}$ chất tăng 1 K. b) S: Đơn vị SI của nhiệt dung riêng là J/kg. c) Đ: Khi giải dạng “Nhận biết khái niệm, đơn vị và ý nghĩa vật lí của nhiệt dung riêng”, cần xác định đúng trạng thái đầu, trạng thái cuối và quy ước dấu hoặc điều kiện áp dụng. d) S: Có thể bỏ qua mọi dữ kiện về khối lượng, nhiệt độ và hiệu suất trong mọi bài toán thuộc dạng này. Kết luận dựa trên định nghĩa, điều kiện áp dụng và quy ước trong SGK KNTT
\item (Sai) Đơn vị SI của nhiệt dung riêng là J/kg.
\item (Đúng) Khi giải dạng “Nhận biết khái niệm, đơn vị và ý nghĩa vật lí của nhiệt dung riêng”, cần xác định đúng trạng thái đầu, trạng thái cuối và quy ước dấu hoặc điều kiện áp dụng.
\item (Sai) Có thể bỏ qua mọi dữ kiện về khối lượng, nhiệt độ và hiệu suất trong mọi bài toán thuộc dạng này.
\end{itemchoice}
}
\end{ex}

% ===== Câu 62 =====
% ID: L12C1B6-159-TL
% Mức: VD
\begin{bt}
% ID: L12C1B6-159-TL
% Mức: VD
Trình bày cách nhận biết và phương pháp giải dạng “Nhận biết khái niệm, đơn vị và ý nghĩa vật lí của nhiệt dung riêng”. Phân tích nhận định sau và sửa lại nếu sai: “Đơn vị SI của nhiệt dung riêng là J/kg.”
\loigiai{
Trước hết xác định hiện tượng, trạng thái và các đại lượng đã cho. Nêu định nghĩa hoặc hệ thức phù hợp, đổi về đơn vị SI, áp dụng đúng điều kiện và kiểm tra ý nghĩa kết quả. Nhận định cần đối chiếu: Nhiệt dung riêng cho biết nhiệt lượng cần để làm $1\,\text{kg}$ chất tăng 1 K. Vì vậy nhận định đã cho sai; phát biểu đúng là: Nhiệt dung riêng cho biết nhiệt lượng cần để làm $1\,\text{kg}$ chất tăng 1 K.
}
\end{bt}

% ===== Câu 65 =====
% ID: L12C1B6-160-DS
% Mức: TH
\begin{ex}
% ID: L12C1B6-160-DS
% Mức: TH
Xét các phát biểu sau về dạng “Nhận biết khái niệm, đơn vị và ý nghĩa vật lí của nhiệt dung riêng” (bộ số 1).
\choiceTF
{Đơn vị SI của nhiệt dung riêng là J/kg.}
{\True Khi giải dạng “Nhận biết khái niệm, đơn vị và ý nghĩa vật lí của nhiệt dung riêng”, cần xác định đúng trạng thái đầu, trạng thái cuối và quy ước dấu hoặc điều kiện áp dụng.}
{Có thể bỏ qua mọi dữ kiện về khối lượng, nhiệt độ và hiệu suất trong mọi bài toán thuộc dạng này.}
{\True Nhiệt dung riêng cho biết nhiệt lượng cần để làm $1\,\text{kg}$ chất tăng 1 K.}
\loigiai{
\begin{itemchoice}
\item (Sai) Đơn vị SI của nhiệt dung riêng là J/kg. (Vì): Đơn vị SI của nhiệt dung riêng là J/kg. b) Đ: Khi giải dạng “Nhận biết khái niệm, đơn vị và ý nghĩa vật lí của nhiệt dung riêng”, cần xác định đúng trạng thái đầu, trạng thái cuối và quy ước dấu hoặc điều kiện áp dụng. c) S: Có thể bỏ qua mọi dữ kiện về khối lượng, nhiệt độ và hiệu suất trong mọi bài toán thuộc dạng này. d) Đ: Nhiệt dung riêng cho biết nhiệt lượng cần để làm $1\,\text{kg}$ chất tăng 1 K. Kết luận dựa trên định nghĩa, điều kiện áp dụng và quy ước trong SGK KNTT
\item (Đúng) Khi giải dạng “Nhận biết khái niệm, đơn vị và ý nghĩa vật lí của nhiệt dung riêng”, cần xác định đúng trạng thái đầu, trạng thái cuối và quy ước dấu hoặc điều kiện áp dụng.
\item (Sai) Có thể bỏ qua mọi dữ kiện về khối lượng, nhiệt độ và hiệu suất trong mọi bài toán thuộc dạng này.
\item (Đúng) Nhiệt dung riêng cho biết nhiệt lượng cần để làm $1\,\text{kg}$ chất tăng 1 K.
\end{itemchoice}
}
\end{ex}

% ===== Câu 66 =====
% ID: L12C1B6-161-TL
% Mức: VDC
\begin{bt}
% ID: L12C1B6-161-TL
% Mức: VDC
Trình bày cách nhận biết và phương pháp giải dạng “Nhận biết khái niệm, đơn vị và ý nghĩa vật lí của nhiệt dung riêng”. Phân tích nhận định sau và sửa lại nếu sai: “Nhiệt dung riêng cho biết nhiệt lượng cần để làm $1\,\text{kg}$ chất tăng 1 K.”
\loigiai{
Trước hết xác định hiện tượng, trạng thái và các đại lượng đã cho. Nêu định nghĩa hoặc hệ thức phù hợp, đổi về đơn vị SI, áp dụng đúng điều kiện và kiểm tra ý nghĩa kết quả. Nhận định cần đối chiếu: Nhiệt dung riêng cho biết nhiệt lượng cần để làm $1\,\text{kg}$ chất tăng 1 K. Vì vậy nhận định đã cho đúng.
}
\end{bt}

% ===== Câu 69 =====
% ID: L12C1B6-162-DS
% Mức: TH
\begin{ex}
% ID: L12C1B6-162-DS
% Mức: TH
Xét các phát biểu sau về dạng “Nhận biết khái niệm, đơn vị và ý nghĩa vật lí của nhiệt dung riêng” (bộ số 5).
\choiceTF
{Đơn vị SI của nhiệt dung riêng là J/kg.}
{\True Khi giải dạng “Nhận biết khái niệm, đơn vị và ý nghĩa vật lí của nhiệt dung riêng”, cần xác định đúng trạng thái đầu, trạng thái cuối và quy ước dấu hoặc điều kiện áp dụng.}
{Có thể bỏ qua mọi dữ kiện về khối lượng, nhiệt độ và hiệu suất trong mọi bài toán thuộc dạng này.}
{\True Nhiệt dung riêng cho biết nhiệt lượng cần để làm $1\,\text{kg}$ chất tăng 1 K.}
\loigiai{
\begin{itemchoice}
\item (Sai) Đơn vị SI của nhiệt dung riêng là J/kg. (Vì): Đơn vị SI của nhiệt dung riêng là J/kg. b) Đ: Khi giải dạng “Nhận biết khái niệm, đơn vị và ý nghĩa vật lí của nhiệt dung riêng”, cần xác định đúng trạng thái đầu, trạng thái cuối và quy ước dấu hoặc điều kiện áp dụng. c) S: Có thể bỏ qua mọi dữ kiện về khối lượng, nhiệt độ và hiệu suất trong mọi bài toán thuộc dạng này. d) Đ: Nhiệt dung riêng cho biết nhiệt lượng cần để làm $1\,\text{kg}$ chất tăng 1 K. Kết luận dựa trên định nghĩa, điều kiện áp dụng và quy ước trong SGK KNTT
\item (Đúng) Khi giải dạng “Nhận biết khái niệm, đơn vị và ý nghĩa vật lí của nhiệt dung riêng”, cần xác định đúng trạng thái đầu, trạng thái cuối và quy ước dấu hoặc điều kiện áp dụng.
\item (Sai) Có thể bỏ qua mọi dữ kiện về khối lượng, nhiệt độ và hiệu suất trong mọi bài toán thuộc dạng này.
\item (Đúng) Nhiệt dung riêng cho biết nhiệt lượng cần để làm $1\,\text{kg}$ chất tăng 1 K.
\end{itemchoice}
}
\end{ex}

% ===== Câu 70 =====
% ID: L12C1B6-163-TL
% Mức: VDC
\begin{bt}
% ID: L12C1B6-163-TL
% Mức: VDC
Trình bày cách nhận biết và phương pháp giải dạng “Nhận biết khái niệm, đơn vị và ý nghĩa vật lí của nhiệt dung riêng”. Phân tích nhận định sau và sửa lại nếu sai: “Nhiệt dung riêng cho biết nhiệt lượng cần để làm $1\,\text{kg}$ chất tăng 1 K.”
\loigiai{
Trước hết xác định hiện tượng, trạng thái và các đại lượng đã cho. Nêu định nghĩa hoặc hệ thức phù hợp, đổi về đơn vị SI, áp dụng đúng điều kiện và kiểm tra ý nghĩa kết quả. Nhận định cần đối chiếu: Nhiệt dung riêng cho biết nhiệt lượng cần để làm $1\,\text{kg}$ chất tăng 1 K. Vì vậy nhận định đã cho đúng.
}
\end{bt}

% ===== Câu 73 =====
% ID: L12C1B6-164-DS
% Mức: VD
\begin{ex}
% ID: L12C1B6-164-DS
% Mức: VD
Xét các phát biểu sau về dạng “Nhận biết khái niệm, đơn vị và ý nghĩa vật lí của nhiệt dung riêng” (bộ số 2).
\choiceTF
{\True Khi giải dạng “Nhận biết khái niệm, đơn vị và ý nghĩa vật lí của nhiệt dung riêng”, cần xác định đúng trạng thái đầu, trạng thái cuối và quy ước dấu hoặc điều kiện áp dụng.}
{Có thể bỏ qua mọi dữ kiện về khối lượng, nhiệt độ và hiệu suất trong mọi bài toán thuộc dạng này.}
{\True Nhiệt dung riêng cho biết nhiệt lượng cần để làm $1\,\text{kg}$ chất tăng 1 K.}
{Đơn vị SI của nhiệt dung riêng là J/kg.}
\loigiai{
\begin{itemchoice}
\item (Đúng) Khi giải dạng “Nhận biết khái niệm, đơn vị và ý nghĩa vật lí của nhiệt dung riêng”, cần xác định đúng trạng thái đầu, trạng thái cuối và quy ước dấu hoặc điều kiện áp dụng. (Vì): Khi giải dạng “Nhận biết khái niệm, đơn vị và ý nghĩa vật lí của nhiệt dung riêng”, cần xác định đúng trạng thái đầu, trạng thái cuối và quy ước dấu hoặc điều kiện áp dụng. b) S: Có thể bỏ qua mọi dữ kiện về khối lượng, nhiệt độ và hiệu suất trong mọi bài toán thuộc dạng này. c) Đ: Nhiệt dung riêng cho biết nhiệt lượng cần để làm $1\,\text{kg}$ chất tăng 1 K. d) S: Đơn vị SI của nhiệt dung riêng là J/kg. Kết luận dựa trên định nghĩa, điều kiện áp dụng và quy ước trong SGK KNTT
\item (Sai) Có thể bỏ qua mọi dữ kiện về khối lượng, nhiệt độ và hiệu suất trong mọi bài toán thuộc dạng này.
\item (Đúng) Nhiệt dung riêng cho biết nhiệt lượng cần để làm $1\,\text{kg}$ chất tăng 1 K.
\item (Sai) Đơn vị SI của nhiệt dung riêng là J/kg.
\end{itemchoice}
}
\end{ex}

% ===== Câu 75 =====
% ID: L12C1B6-165-TLN
% Mức: NB
\begin{ex}
% ID: L12C1B6-165-TLN
% Mức: NB
Cho bốn nhận định: (1) Nhiệt dung riêng cho biết nhiệt lượng cần để làm $1\,\text{kg}$ chất tăng 1 K. (2) Đơn vị SI của nhiệt dung riêng là J/kg. (3) Cần kiểm tra đơn vị và điều kiện áp dụng khi xử lí dạng “Nhận biết khái niệm, đơn vị và ý nghĩa vật lí của nhiệt dung riêng”. (4) Có thể thay tùy ý nhiệt độ Celsius cho nhiệt độ tuyệt đối trong mọi công thức. Có bao nhiêu nhận định đúng? Trả lời bằng một số nguyên.
\shortans{2}
\loigiai{
Nhận định (1) và (3) đúng; (2) và (4) sai. Vậy có 2 nhận định đúng.
}
\end{ex}

% ===== Câu 77 =====
% ID: L12C1B6-166-DS
% Mức: VDC
\begin{ex}
% ID: L12C1B6-166-DS
% Mức: VDC
Xét các phát biểu sau về dạng “Nhận biết khái niệm, đơn vị và ý nghĩa vật lí của nhiệt dung riêng” (bộ số 3).
\choiceTF
{Có thể bỏ qua mọi dữ kiện về khối lượng, nhiệt độ và hiệu suất trong mọi bài toán thuộc dạng này.}
{\True Nhiệt dung riêng cho biết nhiệt lượng cần để làm $1\,\text{kg}$ chất tăng 1 K.}
{Đơn vị SI của nhiệt dung riêng là J/kg.}
{\True Khi giải dạng “Nhận biết khái niệm, đơn vị và ý nghĩa vật lí của nhiệt dung riêng”, cần xác định đúng trạng thái đầu, trạng thái cuối và quy ước dấu hoặc điều kiện áp dụng.}
\loigiai{
\begin{itemchoice}
\item (Sai) Có thể bỏ qua mọi dữ kiện về khối lượng, nhiệt độ và hiệu suất trong mọi bài toán thuộc dạng này. (Vì): Có thể bỏ qua mọi dữ kiện về khối lượng, nhiệt độ và hiệu suất trong mọi bài toán thuộc dạng này. b) Đ: Nhiệt dung riêng cho biết nhiệt lượng cần để làm $1\,\text{kg}$ chất tăng 1 K. c) S: Đơn vị SI của nhiệt dung riêng là J/kg. d) Đ: Khi giải dạng “Nhận biết khái niệm, đơn vị và ý nghĩa vật lí của nhiệt dung riêng”, cần xác định đúng trạng thái đầu, trạng thái cuối và quy ước dấu hoặc điều kiện áp dụng. Kết luận dựa trên định nghĩa, điều kiện áp dụng và quy ước trong SGK KNTT
\item (Đúng) Nhiệt dung riêng cho biết nhiệt lượng cần để làm $1\,\text{kg}$ chất tăng 1 K.
\item (Sai) Đơn vị SI của nhiệt dung riêng là J/kg.
\item (Đúng) Khi giải dạng “Nhận biết khái niệm, đơn vị và ý nghĩa vật lí của nhiệt dung riêng”, cần xác định đúng trạng thái đầu, trạng thái cuối và quy ước dấu hoặc điều kiện áp dụng.
\end{itemchoice}
}
\end{ex}

% ===== Câu 79 =====
% ID: L12C1B6-167-TLN
% Mức: TH
\begin{ex}
% ID: L12C1B6-167-TLN
% Mức: TH
Cho bốn nhận định: (1) Nhiệt dung riêng cho biết nhiệt lượng cần để làm $1\,\text{kg}$ chất tăng 1 K. (2) Đơn vị SI của nhiệt dung riêng là J/kg. (3) Cần kiểm tra đơn vị và điều kiện áp dụng khi xử lí dạng “Nhận biết khái niệm, đơn vị và ý nghĩa vật lí của nhiệt dung riêng”. (4) Có thể thay tùy ý nhiệt độ Celsius cho nhiệt độ tuyệt đối trong mọi công thức. Có bao nhiêu nhận định đúng? Trả lời bằng một số nguyên.
\shortans{2}
\loigiai{
Nhận định (1) và (3) đúng; (2) và (4) sai. Vậy có 2 nhận định đúng.
}
\end{ex}

% ===== Câu 83 =====
% ID: L12C1B6-168-TLN
% Mức: VD
\begin{ex}
% ID: L12C1B6-168-TLN
% Mức: VD
Cho bốn nhận định: (1) Nhiệt dung riêng cho biết nhiệt lượng cần để làm $1\,\text{kg}$ chất tăng 1 K. (2) Đơn vị SI của nhiệt dung riêng là J/kg. (3) Cần kiểm tra đơn vị và điều kiện áp dụng khi xử lí dạng “Nhận biết khái niệm, đơn vị và ý nghĩa vật lí của nhiệt dung riêng”. (4) Có thể thay tùy ý nhiệt độ Celsius cho nhiệt độ tuyệt đối trong mọi công thức. Có bao nhiêu nhận định đúng? Trả lời bằng một số nguyên.
\shortans{2}
\loigiai{
Nhận định (1) và (3) đúng; (2) và (4) sai. Vậy có 2 nhận định đúng.
}
\end{ex}

% ===== Câu 87 =====
% ID: L12C1B6-169-TLN
% Mức: VD
\begin{ex}
% ID: L12C1B6-169-TLN
% Mức: VD
Cho bốn nhận định: (1) Nhiệt dung riêng cho biết nhiệt lượng cần để làm $1\,\text{kg}$ chất tăng 1 K. (2) Đơn vị SI của nhiệt dung riêng là J/kg. (3) Cần kiểm tra đơn vị và điều kiện áp dụng khi xử lí dạng “Nhận biết khái niệm, đơn vị và ý nghĩa vật lí của nhiệt dung riêng”. (4) Có thể thay tùy ý nhiệt độ Celsius cho nhiệt độ tuyệt đối trong mọi công thức. Có bao nhiêu nhận định đúng? Trả lời bằng một số nguyên.
\shortans{2}
\loigiai{
Nhận định (1) và (3) đúng; (2) và (4) sai. Vậy có 2 nhận định đúng.
}
\end{ex}

% ===== Câu 91 =====
% ID: L12C1B6-170-TLN
% Mức: VDC
\begin{ex}
% ID: L12C1B6-170-TLN
% Mức: VDC
Cho bốn nhận định: (1) Nhiệt dung riêng cho biết nhiệt lượng cần để làm $1\,\text{kg}$ chất tăng 1 K. (2) Đơn vị SI của nhiệt dung riêng là J/kg. (3) Cần kiểm tra đơn vị và điều kiện áp dụng khi xử lí dạng “Nhận biết khái niệm, đơn vị và ý nghĩa vật lí của nhiệt dung riêng”. (4) Có thể thay tùy ý nhiệt độ Celsius cho nhiệt độ tuyệt đối trong mọi công thức. Có bao nhiêu nhận định đúng? Trả lời bằng một số nguyên.
\shortans{2}
\loigiai{
Nhận định (1) và (3) đúng; (2) và (4) sai. Vậy có 2 nhận định đúng.
}
\end{ex}

% ===== Câu 95 =====
% ID: L12C1B6-171-TLN
% Mức: VDC
\begin{ex}
% ID: L12C1B6-171-TLN
% Mức: VDC
Cho bốn nhận định: (1) Nhiệt dung riêng cho biết nhiệt lượng cần để làm $1\,\text{kg}$ chất tăng 1 K. (2) Đơn vị SI của nhiệt dung riêng là J/kg. (3) Cần kiểm tra đơn vị và điều kiện áp dụng khi xử lí dạng “Nhận biết khái niệm, đơn vị và ý nghĩa vật lí của nhiệt dung riêng”. (4) Có thể thay tùy ý nhiệt độ Celsius cho nhiệt độ tuyệt đối trong mọi công thức. Có bao nhiêu nhận định đúng? Trả lời bằng một số nguyên.
\shortans{2}
\loigiai{
Nhận định (1) và (3) đúng; (2) và (4) sai. Vậy có 2 nhận định đúng.
}
\end{ex}

% ===== Câu 132 =====
% ID: L12C1B6-172-TN
% Mức: NB
\begin{ex}
% ID: L12C1B6-172-TN
% Mức: NB
Câu 1 thuộc dạng “Nhận biết khái niệm, đơn vị và ý nghĩa vật lí của nhiệt dung riêng”. Phát biểu nào sau đây đúng?
\choice
{Đơn vị SI của nhiệt dung riêng là J/kg.}
{Nhiệt lượng và nhiệt độ là hai đại lượng hoàn toàn giống nhau.}
{Mọi quá trình nhiệt học đều không phụ thuộc điều kiện thực hiện.}
{\True Nhiệt dung riêng cho biết nhiệt lượng cần để làm $1\,\text{kg}$ chất tăng 1 K.}
\loigiai{
Phát biểu đúng là: Nhiệt dung riêng cho biết nhiệt lượng cần để làm $1\,\text{kg}$ chất tăng 1 K. Các phương án còn lại trái với mô hình, định nghĩa hoặc hệ thức nhiệt học tương ứng.
}
\end{ex}

% ===== Câu 136 =====
% ID: L12C1B6-173-TN
% Mức: NB
\begin{ex}
% ID: L12C1B6-173-TN
% Mức: NB
Câu 5 thuộc dạng “Nhận biết khái niệm, đơn vị và ý nghĩa vật lí của nhiệt dung riêng”. Phát biểu nào sau đây đúng?
\choice
{Đơn vị SI của nhiệt dung riêng là J/kg.}
{Nhiệt lượng và nhiệt độ là hai đại lượng hoàn toàn giống nhau.}
{Mọi quá trình nhiệt học đều không phụ thuộc điều kiện thực hiện.}
{\True Nhiệt dung riêng cho biết nhiệt lượng cần để làm $1\,\text{kg}$ chất tăng 1 K.}
\loigiai{
Phát biểu đúng là: Nhiệt dung riêng cho biết nhiệt lượng cần để làm $1\,\text{kg}$ chất tăng 1 K. Các phương án còn lại trái với mô hình, định nghĩa hoặc hệ thức nhiệt học tương ứng.
}
\end{ex}

% ===== Câu 140 =====
% ID: L12C1B6-174-TN
% Mức: NB
\begin{ex}
% ID: L12C1B6-174-TN
% Mức: NB
Câu 9 thuộc dạng “Nhận biết khái niệm, đơn vị và ý nghĩa vật lí của nhiệt dung riêng”. Phát biểu nào sau đây đúng?
\choice
{Đơn vị SI của nhiệt dung riêng là J/kg.}
{Nhiệt lượng và nhiệt độ là hai đại lượng hoàn toàn giống nhau.}
{Mọi quá trình nhiệt học đều không phụ thuộc điều kiện thực hiện.}
{\True Nhiệt dung riêng cho biết nhiệt lượng cần để làm $1\,\text{kg}$ chất tăng 1 K.}
\loigiai{
Phát biểu đúng là: Nhiệt dung riêng cho biết nhiệt lượng cần để làm $1\,\text{kg}$ chất tăng 1 K. Các phương án còn lại trái với mô hình, định nghĩa hoặc hệ thức nhiệt học tương ứng.
}
\end{ex}

% ===== Câu 143 =====
% ID: L12C1B6-175-TN
% Mức: TH
\begin{ex}
% ID: L12C1B6-175-TN
% Mức: TH
Câu 2 thuộc dạng “Nhận biết khái niệm, đơn vị và ý nghĩa vật lí của nhiệt dung riêng”. Phát biểu nào sau đây đúng?
\choice
{Nhiệt lượng và nhiệt độ là hai đại lượng hoàn toàn giống nhau.}
{Mọi quá trình nhiệt học đều không phụ thuộc điều kiện thực hiện.}
{\True Nhiệt dung riêng cho biết nhiệt lượng cần để làm $1\,\text{kg}$ chất tăng 1 K.}
{Đơn vị SI của nhiệt dung riêng là J/kg.}
\loigiai{
Phát biểu đúng là: Nhiệt dung riêng cho biết nhiệt lượng cần để làm $1\,\text{kg}$ chất tăng 1 K. Các phương án còn lại trái với mô hình, định nghĩa hoặc hệ thức nhiệt học tương ứng.
}
\end{ex}

% ===== Câu 145 =====
% ID: L12C1B6-176-TN
% Mức: TH
\begin{ex}
% ID: L12C1B6-176-TN
% Mức: TH
Câu 6 thuộc dạng “Nhận biết khái niệm, đơn vị và ý nghĩa vật lí của nhiệt dung riêng”. Phát biểu nào sau đây đúng?
\choice
{Nhiệt lượng và nhiệt độ là hai đại lượng hoàn toàn giống nhau.}
{Mọi quá trình nhiệt học đều không phụ thuộc điều kiện thực hiện.}
{\True Nhiệt dung riêng cho biết nhiệt lượng cần để làm $1\,\text{kg}$ chất tăng 1 K.}
{Đơn vị SI của nhiệt dung riêng là J/kg.}
\loigiai{
Phát biểu đúng là: Nhiệt dung riêng cho biết nhiệt lượng cần để làm $1\,\text{kg}$ chất tăng 1 K. Các phương án còn lại trái với mô hình, định nghĩa hoặc hệ thức nhiệt học tương ứng.
}
\end{ex}

% ===== Câu 147 =====
% ID: L12C1B6-177-TN
% Mức: TH
\begin{ex}
% ID: L12C1B6-177-TN
% Mức: TH
Câu 10 thuộc dạng “Nhận biết khái niệm, đơn vị và ý nghĩa vật lí của nhiệt dung riêng”. Phát biểu nào sau đây đúng?
\choice
{Nhiệt lượng và nhiệt độ là hai đại lượng hoàn toàn giống nhau.}
{Mọi quá trình nhiệt học đều không phụ thuộc điều kiện thực hiện.}
{\True Nhiệt dung riêng cho biết nhiệt lượng cần để làm $1\,\text{kg}$ chất tăng 1 K.}
{Đơn vị SI của nhiệt dung riêng là J/kg.}
\loigiai{
Phát biểu đúng là: Nhiệt dung riêng cho biết nhiệt lượng cần để làm $1\,\text{kg}$ chất tăng 1 K. Các phương án còn lại trái với mô hình, định nghĩa hoặc hệ thức nhiệt học tương ứng.
}
\end{ex}

% ===== Câu 149 =====
% ID: L12C1B6-178-TN
% Mức: VD
\begin{ex}
% ID: L12C1B6-178-TN
% Mức: VD
Câu 3 thuộc dạng “Nhận biết khái niệm, đơn vị và ý nghĩa vật lí của nhiệt dung riêng”. Phát biểu nào sau đây đúng?
\choice
{Mọi quá trình nhiệt học đều không phụ thuộc điều kiện thực hiện.}
{\True Nhiệt dung riêng cho biết nhiệt lượng cần để làm $1\,\text{kg}$ chất tăng 1 K.}
{Đơn vị SI của nhiệt dung riêng là J/kg.}
{Nhiệt lượng và nhiệt độ là hai đại lượng hoàn toàn giống nhau.}
\loigiai{
Phát biểu đúng là: Nhiệt dung riêng cho biết nhiệt lượng cần để làm $1\,\text{kg}$ chất tăng 1 K. Các phương án còn lại trái với mô hình, định nghĩa hoặc hệ thức nhiệt học tương ứng.
}
\end{ex}

% ===== Câu 151 =====
% ID: L12C1B6-179-TN
% Mức: VD
\begin{ex}
% ID: L12C1B6-179-TN
% Mức: VD
Câu 7 thuộc dạng “Nhận biết khái niệm, đơn vị và ý nghĩa vật lí của nhiệt dung riêng”. Phát biểu nào sau đây đúng?
\choice
{Mọi quá trình nhiệt học đều không phụ thuộc điều kiện thực hiện.}
{\True Nhiệt dung riêng cho biết nhiệt lượng cần để làm $1\,\text{kg}$ chất tăng 1 K.}
{Đơn vị SI của nhiệt dung riêng là J/kg.}
{Nhiệt lượng và nhiệt độ là hai đại lượng hoàn toàn giống nhau.}
\loigiai{
Phát biểu đúng là: Nhiệt dung riêng cho biết nhiệt lượng cần để làm $1\,\text{kg}$ chất tăng 1 K. Các phương án còn lại trái với mô hình, định nghĩa hoặc hệ thức nhiệt học tương ứng.
}
\end{ex}

% ===== Câu 153 =====
% ID: L12C1B6-180-TN
% Mức: VDC
\begin{ex}
% ID: L12C1B6-180-TN
% Mức: VDC
Câu 4 thuộc dạng “Nhận biết khái niệm, đơn vị và ý nghĩa vật lí của nhiệt dung riêng”. Phát biểu nào sau đây đúng?
\choice
{\True Nhiệt dung riêng cho biết nhiệt lượng cần để làm $1\,\text{kg}$ chất tăng 1 K.}
{Đơn vị SI của nhiệt dung riêng là J/kg.}
{Nhiệt lượng và nhiệt độ là hai đại lượng hoàn toàn giống nhau.}
{Mọi quá trình nhiệt học đều không phụ thuộc điều kiện thực hiện.}
\loigiai{
Phát biểu đúng là: Nhiệt dung riêng cho biết nhiệt lượng cần để làm $1\,\text{kg}$ chất tăng 1 K. Các phương án còn lại trái với mô hình, định nghĩa hoặc hệ thức nhiệt học tương ứng.
}
\end{ex}

% ===== Câu 155 =====
% ID: L12C1B6-181-TN
% Mức: VDC
\begin{ex}
% ID: L12C1B6-181-TN
% Mức: VDC
Câu 8 thuộc dạng “Nhận biết khái niệm, đơn vị và ý nghĩa vật lí của nhiệt dung riêng”. Phát biểu nào sau đây đúng?
\choice
{\True Nhiệt dung riêng cho biết nhiệt lượng cần để làm $1\,\text{kg}$ chất tăng 1 K.}
{Đơn vị SI của nhiệt dung riêng là J/kg.}
{Nhiệt lượng và nhiệt độ là hai đại lượng hoàn toàn giống nhau.}
{Mọi quá trình nhiệt học đều không phụ thuộc điều kiện thực hiện.}
\loigiai{
Phát biểu đúng là: Nhiệt dung riêng cho biết nhiệt lượng cần để làm $1\,\text{kg}$ chất tăng 1 K. Các phương án còn lại trái với mô hình, định nghĩa hoặc hệ thức nhiệt học tương ứng.
}
\end{ex}
